\begin{statement}{A.26}{Proved}{All odd-prime stable packets lose the proper first layer}
Let $\ell$ be odd and let $s_b$ be any original cusp
packet section in Part A. On the stable $S_3$ cover
$C:v^2=w^7-w$ and its three genus-two reflection
quotients, restrict to any open subgroup $\Delta_0$
over which the selected cusp lifts and packet
actions are defined. Then
\[
 \operatorname{im}\!\left(
    P_{D,\ell}^{s_b(\Delta_0)}
       \longrightarrow H_1(D_{\overline K},\Z_\ell)
                    \right)=0,                         \tag{A.26.1}
\]
and
\[
 \operatorname{im}\!\left(
    P_{C,\ell}^{s_b(\Delta_0)}
       \longrightarrow H_1(C\setminus\{P_0,P_\infty\},
                           \Z_\ell)
                    \right)
              =\Z_\ell[c_C].                           \tag{A.26.2}
\]
The first formula refers to each of the three
reflection quotients. Thus a fixed element on $C$
can be multiplied by an actual fixed cusp power
to put it in $\gamma_2P_{C,\ell}$; a fixed element
on $D$ already lies in $\gamma_2P_{D,\ell}$.
These conclusions include every exceptional packet
left by the determinant tests in A.24.
They do not identify the full nonabelian fixed
subgroup in those exceptional cases.
\end{statement}
\begin{proof}
It suffices first to treat the model $D$ of A.25.
Let $u$ be fixed and let $v$ be its rational proper
homology class. Under the degree-three punctured
map $f:Y\to E\setminus O$, its image is fixed by
the restriction of the original $s_b$.
The input R1, with its finite-index extension
proved in A.22, puts that image in the selected
elliptic cusp inertia. Consequently $f_*v=0$.
The integral conclusion will follow from the
rational one because proper homology is
torsion-free over $\Z_\ell$.

Use the stable reduction above two proved in A.25.
If $v\ne0$, the component degrees $(1,2)$ force
its projections to both elliptic components to
be nonzero. Unlike the normalizer proof, we must
now exhibit the needed monodromy inside
$\Delta_0$, not merely inside an open arithmetic
group. Choose a finite number field extension
whose intersection with $\Delta$ is contained
in $\Delta_0$, and enlarge its completion at two
until the stable model and its components are
defined. This does not change the two positive
component degrees.

Let $\tau$ be a tame $\ell$-inertia element with
nonzero monodromy and $F$ a Frobenius element in
this local Galois group. On the semistable
prime-to-two outer fundamental group, the
commutator $[F,\tau]$ has monodromy equal to
$(q-1)$ times that of $\tau$, where $q>1$ is the
residue cardinality. This follows either from
the tame relation $F\tau F^{-1}=\tau^q$ modulo
wild inertia, or from the edge-twist description
in \cite[Definition 3.1.5 and Theorem 3.1.6]{fghBettsDogra}.
Wild inertia has trivial semistable outer
prime-to-two action. The scalar $q-1$ is nonzero
in $\Q_\ell$, whether or not it is an $\ell$-adic
unit.

The local commutator lies in the global closed
commutator subgroup and in the chosen finite
extension subgroup; hence it is in $\Delta_0$.
The fixedness of $u$ makes its image in the
proper rational unipotent group invariant up
to conjugacy for this monodromy. Changing from
the original cusp base point to the smooth
component base point in A.25 changes the
representative by an inner automorphism only.
The degree-two and degree-three comparison
in that proof therefore gives the same
necessary condition
\[
                 T(v)\in[v,\operatorname{gr}_2
                                      \mathfrak p_D].
\]
The explicit map (A.25.10) kills the right side
and maps the left to the nonzero scalar multiple
$(q-1)m[[x,y],x]$. This contradiction proves
$v=0$, and then (A.26.1).

The other two reflection quotients are obtained
by conjugating the reflection by the cubic
rotation of $C$. After one common finite field
extension the same argument applies to their
induced packet actions and the original map
to $E\setminus O$.
The three pushforward maps from proper
$H_1(C,\Q_\ell)$ to these quotient homologies
are jointly injective. Indeed A.4 gives the
$S_3$ representation as a two-dimensional
trivial part and a standard two-dimensional
representation tensored with a two-dimensional
elliptic homology space. The quotient by a
reflection detects its positive eigenspace.
The positive reflection eigenspaces span the
standard representation, and the trivial part
is detected by every quotient. Equivalently
the sum of the three operators $1+j$ is
invertible on this representation.

The proper homology image of a fixed element
of $P_{C,\ell}$ therefore vanishes. Its open
homology class lies in the boundary lattice,
which is precisely the primitive rank-one
lattice $\Z_\ell[c_C]$. Every cusp power is
fixed by $\Delta_0$, so equality holds in
(A.26.2), and subtraction of that actual
integral cusp power has the claimed effect
on the lower-central filtration.
\end{proof}
